% problem-000216 9 取代基效应与反应速率
\section*{9. 取代基效应与反应速率}
在上个世纪五⼗年代，Corey课题组在全合成�-檀⾹烯(�-Santalene)中⾸先将右旋樟脑(1)的�位溴化，⽣成�-溴樟脑(2)；接着在氯磺酸中溴化转化为3，再利⽤锌粉选择性还原碳基�位溴、最终形成(+)-4。这种制备光学纯化合物(+)-4 的迂回策略是必须的，因为右旋樟脑(1)直接在氯磺酸中溴化则会得到部分外消旋化的 4。

（图：）

\subsection*{9-1}
画出从2⽣成3的关键中间体。

\subsection*{9-2}
画出从(+)-1 ⽣成(−)-4 的关键中间体。

\subsection*{9-3}
简要解释为何化合物2在溴化重排过程中没有发⽣外消旋化；⽽(+)-1则发⽣了外消旋化。
